\documentclass[./测控三级项目.tex]{subfiles} % 使用 subfiles 文档类，指定主文档

\begin{document}
	\subsection{题目重述}
		\begin{compactenum}
			\item 选择适当参数设计一个$y=-(x_1+x_2+10x_3)$的三路加法器电路（忽略负号），完成Multisim仿真测试。
			\item 选择适当元件焊装上述$y=-(x_1+x_2+10x_3)$的三路加法器电路，在实际电路中完成要求的测试。
		\end{compactenum}
	\subsection{题目分析与实施方案}
		\subsubsection{第一问}
			第一问需要构建一个三路加法器电路，并且实现$y=-(x_1+x_2+10x_3)$的功能。然后对这个电路在不同输入信号下的放大效果进行测试，输入信号包括不同频率的正弦波加和，通过仿真来观察输出信号的特性。
		\subsubsection{第二问}
			第二问主要是对焊装后的三路加法器电路进行多方面测试，并对比与仿真波形的差异。要求在实际电路中测试对不同的输入信号进行计算，要求能掌握焊接后的电路情况，依据之前设计的三路加法器电路原理图，在PCB上进行电路焊装。对不同峰值和频率的正弦波、矩形波、三角波信号放大测试，能全面了解电路在实际焊装后的性能表现，包括频率响应、压摆率限制以及阶跃响应等特性，进而分析实际电路与理想仿真情况的偏差原因。
	
	\subsection{理论分析与仿真}
		\subsubsection{第一问}
			图\ref{fig:三路加法器电路}是一个基本的三路加法器电路。
			\begin{figure}[h]
				\centering
				\includegraphics[width=0.45\textwidth]{三路加法器电路}
				\caption{三路加法器电路原理图}
				\label{fig:三路加法器电路}
			\end{figure}\par
			由于差分输入电阻$r_{id}$较大，一般认为同相和反相输入端直接没有电流流经，直接忽略该电流。于是，可依据电路结构，写出同相输入端$v_P$和反相输入端$v_N$的电压公式。
			\begin{align}
				v_P&=0 \\
				\frac{v_1-v_N}{R_1}+\frac{v_2-v_N}{R_2}+\frac{v_3-v_N}{R_3}&=\frac{v_N-v_O}{R_f}
			\end{align}\par
			结合等式$v_P=v_N$，联立可得到
			\begin{align}
				v_O=-(\frac{R_f}{R_1}v_1+\frac{R_f}{R_2}v_2+\frac{R_f}{R_3}v_3)
			\end{align}\par
			要实现题目要求中的$y=-(x_1+x_2+10x_3)$，则需要满足等式$R_f=R_1=R_2=10R_3$即可。\par
			以上理论完备，在Multisim软件中进行仿真实验。
			根据图\ref{fig:三路加法器电路}，搭建仿真电路图，仿真电路使用OP07芯片作为运放芯片，采用$\pm15$V正负电压作为电源供电，反馈电阻网络选择$R_f=R_1=R_2=10$k$\Omega$，$R_3=1$k$\Omega$电阻将四通道信号发生器接入同相放大电路输入端，示波器通道A、B、C接入三路加法器的三路输入端，示波器通道D接入三路加法器输出端，其搭建示意图如图\ref{fig:三路加法器电路仿真原理图}所示。\par
			\begin{figure}[h]
				\centering
				\includegraphics[width=0.45\textwidth]{./6.第六题示波器现象/三路加法器电路仿真原理图}
				\caption{三路加法器电路仿真原理图}
				\label{fig:三路加法器电路仿真原理图}
			\end{figure}\par
			该仿真电路图测得对应的仿真直流和正弦波形如图\ref{fig:仿真三路加法器运算输入输出信号}所示。
			\begin{figure}[H]
				\centering
				\begin{subfigure}{0.4\linewidth}
					\centering
					\includegraphics[width=0.75\linewidth]{./6.第六题示波器现象/仿真第一问}
					\caption{仿真3路50mV直流信号加和}
					\label{fig:仿真3路50mV直流信号加和}
				\end{subfigure}
				\begin{subfigure}{0.4\linewidth}
					\centering
					\includegraphics[width=0.75\linewidth]{./6.第六题示波器现象/仿真第二问}
					\caption{仿真20mV、10Hz正弦信号与20mV、10kHz正弦信号加和}
					\label{fig:仿真20mV、10Hz正弦信号与20mV、10kHz正弦信号加和}
				\end{subfigure}
				\begin{subfigure}{0.4\linewidth}
					\centering
					\includegraphics[width=0.75\linewidth]{./6.第六题示波器现象/仿真第三问}
					\caption{仿真20mV、10Hz正弦信号与100mV、10Hz正弦信号加和}
					\label{fig:仿真20mV、10Hz正弦信号与100mV、10Hz正弦信号加和}
				\end{subfigure}
				\begin{subfigure}{0.4\linewidth}
					\centering
					\includegraphics[width=0.75\linewidth]{./6.第六题示波器现象/仿真第四问}
					\caption{仿真30mV、5kHz正弦信号与200mV、2kHz正弦信号加和}
					\label{fig:仿真30mV、5kHz正弦信号与200mV、2kHz正弦信号加和}
				\end{subfigure}
			\end{figure}
			\begin{figure}[H]\ContinuedFloat
				\centering
				\begin{subfigure}{0.4\linewidth}
					\centering
					\includegraphics[width=0.75\linewidth]{./6.第六题示波器现象/仿真第五问}
					\caption{仿真50mV、1kHz正弦信号、200mV、200Hz正弦信号和十倍20mV、500Hz正弦信号加和}
					\label{仿真50mV、1kHz正弦信号、200mV、200Hz正弦信号和十倍20mV、500Hz正弦信号加和}
				\end{subfigure}
				\caption{仿真三路加法器运算输入输出信号}
				\label{fig:仿真三路加法器运算输入输出信号}
			\end{figure}\par
			对于四通道示波器，波形会比较凌乱可以根据通道对应测量电压之间的运算关系判断电路正确性。若定义$i\in\{A,B,C,D\}$通道的电压为$v_i$，则图示电路的测量电压均基本满足等式$-10v_A-v_B-v_C=v_D$。因此满足题目要求，同时因为输入电压均幅值较小，频率较低，因此不受到运放压摆率的限制造成信号失真。
	\subsection{实际电路测试结果与数据分析}
		\subsubsection{第二问}
			根据等式$R_f=R_1=R_2=10R_3$，若其中$R_3$选用阻值为$=1$k$\Omega$的电阻，则可以得到$R_f$、$R_1$和$R_2$的阻值都为$1$k$\Omega$，与标准电阻值表比对后，实际选用电阻阻值为$R=1$k$\Omega$(Standeard Value,$1\%$)的电阻。\par
			实际电路板除了以上$R_1$、$R_2$、$R_3$、$R_f$电阻以外还存在输入平衡电阻$R$，该电阻是为了维持运算放大器输入级差分电路的对称性，因此该电阻等于负输入端的对地等效电阻，即$R=\frac{1}{\frac{1}{R_1}+\frac{1}{R_2}+\frac{1}{R_3}+\frac{1}{R_f}}=\frac{10}{13}$k$\Omega$。与标准电阻值表比对后，实际选用电阻阻值为$R=768\Omega$(Standeard Value,$1\%$)的电阻。\par
			焊接后，经检查分立元件位置端正、放置平整、焊点表面光滑、无凸起毛刺、无虚焊，周围无助焊剂残留；多引脚器件无短路，焊点均匀一致。符合测试要求。\par
			将信号发生器的三路信号分别加载在三路加法器的信号输入端，将示波器通道CH1接入加法器的输出端。添加FFT算子计算输出通道CH1的幅频曲线。\par
			将电路连接完毕后，先打开线性电源输出开关，随后根据《测控仪器电路项目作业指导书》中要求，加载对应直流信号、正弦信号。\par
			三路加法器实验实测对应的仿真直流信号、正弦波形如图\ref{fig:实际三路加法器运算输入输出信号}所示。
			\begin{figure}[H]
				\centering
				\begin{subfigure}{0.4\linewidth}
					\centering
					\includegraphics[width=0.75\linewidth]{./6.第六题示波器现象/UpoPrtsc-001}
					\caption{实际3路50mV直流信号加和}
					\label{fig:实际3路50mV直流信号加和}
				\end{subfigure}
				\begin{subfigure}{0.4\linewidth}
					\centering
					\includegraphics[width=0.75\linewidth]{./6.第六题示波器现象/UpoPrtsc-002}
					\caption{实际20mV、10Hz正弦信号与20mV、10kHz正弦信号加和}
					\label{fig:实际20mV、10Hz正弦信号与20mV、10kHz正弦信号加和}
				\end{subfigure}
			\end{figure}
			\begin{figure}[H]\ContinuedFloat
			\centering
				\begin{subfigure}{0.4\linewidth}
					\centering
					\includegraphics[width=0.75\linewidth]{./6.第六题示波器现象/UpoPrtsc-003}
					\caption{实际20mV、10Hz正弦信号与100mV、10Hz正弦信号加和}
					\label{fig:实际20mV、10Hz正弦信号与100mV、10Hz正弦信号加和}
				\end{subfigure}
				\begin{subfigure}{0.4\linewidth}
					\centering
					\includegraphics[width=0.75\linewidth]{./6.第六题示波器现象/UpoPrtsc-004}
					\caption{实际30mV、5kHz正弦信号与200mV、2kHz正弦信号加和}
					\label{fig:实际30mV、5kHz正弦信号与200mV、2kHz正弦信号加和}
				\end{subfigure}
				\begin{subfigure}{0.4\linewidth}
					\centering
					\includegraphics[width=0.75\linewidth]{./6.第六题示波器现象/UpoPrtsc-005}
					\caption{实际50mV、1kHz正弦信号、200mV、200Hz正弦信号和十倍20mV、500Hz正弦信号加和}
					\label{实际50mV、1kHz正弦信号、200mV、200Hz正弦信号和十倍20mV、500Hz正弦信号加和}
				\end{subfigure}
				\caption{实际三路加法器运算输入输出信号}
				\label{fig:实际三路加法器运算输入输出信号}
			\end{figure}\par
			通过FFT变换，可以将时域信号转换为频域信号，观察不同频率信号下的强度，以此来检查加法器是否处于正常的工作状态。
	\subsection{结论与收获}
		\subsubsection{结论}
		\paragraph{\qquad 电路功能实现结论}
		通过 Multisim 仿真和实际电路测试，所设计的三路加法器电路能够实现对不同类型信号的加法运算。对于直流信号，如 3 个 50mV 直流信号加和，在合理的误差范围内得到了接近理论值的结果，验证了电路对直流信号相加的准确性。\par
		在正弦信号相加方面，无论是不同频率还是同频率正弦信号的组合，电路均能输出相应的叠加波形，其频率成分和幅度变化趋势与理论分析基本相符，表明电路对正弦信号的加法功能正常。例如，1 个 20mV、10Hz 正弦信号与 1 个 100mV、10Hz 正弦信号单倍加和后得到接近 120mV、10Hz 的正弦波，1 个 30mV、5KHz 正弦信号与 1 个 200mV、2kHz 正弦信号单倍加和后输出复杂周期性波形，其幅度和周期特性符合理论预期。\par
		对于多个不同频率正弦信号放大 10 倍后的加和，如 1 个 50mV、1KHz 正弦信号、1 个 200mV、200Hz 正弦信号和 1 个 20mV、500Hz 正弦信号的情况，电路也能产生包含相应频率成分且幅度在合理范围内变化的复杂波形，实现了多信号的加法运算功能。
		\paragraph{\qquad 误差分析结论}
		在实际电路测试中，与 Multisim 仿真结果相比存在一定误差。对于直流信号加和，误差主要来源于元件参数误差（如电阻实际值与标称值偏差）、运放失调电压以及电源稳定性等因素。这些因素导致实际输出直流电压值与理论值有偏差，但在可接受范围内，说明电路设计基本合理，但仍有优化空间。\par
		在正弦信号相加测试中，误差原因较为复杂。一方面，电路的频率响应特性对不同频率信号的增益和相位影响不同，导致实际输出波形在幅度和相位上与理论值存在差异。例如，高频信号（如 10KHz 正弦信号）在通过电路时可能因电容、电感等元件的影响而产生衰减和相移，使叠加后的波形与理论计算不完全一致。另一方面，运放的非理想特性，如带宽限制、非线性特性等也会对输出波形产生影响。带宽有限会使高频信号放大倍数下降，非线性特性会导致信号失真，尤其是在信号幅度较大时更为明显。此外，实际电路中的寄生参数（如布线寄生电容、电感）也会干扰信号传输，影响波形质量。
		\subsubsection{收获}
		\paragraph{\qquad 理论分析层面}
		小组成员深入理解了电路加法运算的原理，掌握了不同类型信号（直流、不同频率正弦信号）相加的理论计算方法，能够准确推导输出信号的预期结果，如不同频率正弦信号叠加后的波形特征、幅度变化范围等。通过分析实际电路与仿真结果的误差，对电路中的各种非理想因素（如元件参数误差、运放特性、频率响应等）有了更深刻的认识，明白了这些因素如何影响电路性能，学会了从理论角度解释和评估电路的实际表现，为今后电路设计的优化提供了理论依据。同时，在处理多个信号的加法运算过程中，进一步巩固了信号与系统的相关知识，如信号频率、相位、幅度之间的关系以及它们在电路中的变化规律。
		\paragraph{\qquad 动手实践方面}
		在电路设计阶段，学会根据需求选择合适的元件参数，设计出满足功能要求的三路加法器电路，提高了电路设计能力。在焊装电路过程中，熟练掌握了电路焊接技巧，能够正确安装和连接各种元件，确保电路的物理连接正确可靠，培养了实际动手操作能力。使用示波器进行测试时，学会了设置示波器的各种参数（如时间基、幅度、触发方式等），准确测量和观察信号波形，熟练运用示波器的存储功能导出测试波形，便于后续分析。通过对实际电路的调试和优化，积累了丰富的电路故障排查和解决问题的经验，如当发现输出波形异常时，能够从元件焊接、参数选择、电路布局等方面进行检查和调整，提高了工程实践能力。
		\paragraph{\qquad 精神层面}
		在整个项目过程中，小组成员培养了团队合作精神，大家分工协作，共同完成电路设计、仿真、焊装、测试等任务。在遇到问题时，能够积极沟通交流，共同探讨解决方案，提高了团队的凝聚力和协作效率。面对实际电路测试中出现的各种误差和问题，成员们没有退缩，而是坚持不懈地进行分析和调试，培养了勇于面对困难和挑战的精神。通过不断尝试和改进，最终成功实现电路功能并达到较好的性能指标，增强了自信心和成就感，激发了对电子电路设计的兴趣和热情。同时，严谨细致的实验态度贯穿始终，从电路设计计算到元件选择、焊接以及测试过程中的每一个环节，都要求成员们认真对待，避免因粗心大意而导致错误，这种态度将对今后的学习和工作产生积极影响
	
	
\end{document}